% Review contribution: ordinary realization and inherited bi-interpretation.
% Not a new arithmetic-bootstrap proof. Citation keys may be adapted.
\section{Ordinary realization and the precise iPRA equivalence}

The new result changes one introduction selector.  Its arithmetic bridge is
the existing ConciseIPRA bridge, and is transferred by equality of the
\emph{entire contextual consequence relation}.  This distinction prevents
agreement on the standard natural numbers from being mistaken for
bi-interpretability.

\subsection{A completely specified ordinary realization}
Elaborate iterator syntax into a staged first-order signature.  Start with
$0,S$.  For each fixed earlier-stage term $h(u,\vec p)$ introduce a symbol
$i_h(n,a,\vec p)$.  Its defining axioms are the universal closures of
\[
 i_h(0,a,\vec p)=a,\qquad
 i_h(Sn,a,\vec p)=h(i_h(n,a,\vec p),\vec p).
\]
Only finite descriptions are used.  There is no variable ranging over
function descriptions and no universally total evaluator.

Let $T_I$ have ordinary intuitionistic first-order logic with equality,
the displayed computation schemas, the two axioms
\[
 \forall x\bigl(x=0\lor\exists y\,x=Sy\bigr),\qquad
 (0=S0)\to\bot,
\]
and contextual literal-equation induction
\[
 \frac{\Gamma\vdash E(0)\qquad\Gamma,E(z)\vdash E(Sz)}
      {\Gamma\vdash E(t)}.
\]
Here $E$ is a literal equation and $z$ is fresh for the retained context,
the target, and the passive parameters.  No induction on a quantified
formula is primitive.

For completeness, the ordinary logic consists of reusable hypotheses;
$\land$ introduction and both projections; both $\lor$ introductions and
two-branch elimination; implication introduction discharging one
assumption and modus ponens; universal introduction at a fresh
eigenparameter and universal specialization at any term; existential
introduction at a term and existential elimination at a fresh
eigenparameter; falsity elimination; equality reflexivity; and equality
replacement in any formula.  Context weakening, exchange, contraction,
capture-free substitution, and cut are admissible.  Freshness for
universal introduction concerns the retained context; freshness for
existential elimination concerns both the retained context and its
conclusion.  These are all the logical and nonlogical primitives of
$T_I$.

The ordinary reading of a native formula is
\[
 (s=t)^\circ=(s=t),\qquad
 (Q_yB)^\circ=B^\circ(0)\to\forall y\,B^\circ(y),
\]
\[
 (D_{x,y}B)^\circ=
 \exists x\bigl(B^\circ(x,0)\to\forall y\,B^\circ(x,y)\bigr).
\]
The interpretation domain is $x=x$ and equality is identity.  The
designated native contradiction $0=S0$ interprets as an actual
contradiction because the separation axiom is explicit.  Omitting that
axiom would leave the ordinary positive arithmetic axioms compatible
with a singleton structure.

The new mirrored introduction is sound in $T_I$: in its progressive
branch its matrix is a literal equation, so contextual equation
induction applies under the assumed base.  In its case branch numerical
coverage combines the zero case and all successor cases.  The other
five rule families have their displayed equality or ordinary logical
readings.

\subsection{Explicit inverse logical macros}\label{sec:choice-introduction}
One can first erase auxiliary $Q$ occurrences by replacing
$Q_yB(y)$ with $D_{u,y}B(y)$, with fresh unused $u$.  These formulas are
interderivable: case introduction gives the forward implication, and
paired elimination returns the opened closure in the reverse direction.
Congruence under $Q$ and $D$ then gives this replacement inside every
native formula.  Thus the following macros need only the Eq/$D$
fragment, without restricting the public language.

Write $C(v,a,b)=I(v,a;u.b)$, with $u$ unused, and $1=S0$.  On Eq/$D$
formulas define the gate
\[
 G_v(s=t):=\bigl(C(v,0,s)=C(v,0,t)\bigr),\qquad
 G_v(D_{x,y}A):=D_{x,y}G_vA,
\]
renaming all binders away from $v$.  Then $G_0A$ is a theorem,
$G_{Sv}A\dashv\vdash A$, and
$A\dashv\vdash D_{u,y}G_yA$ for fresh unused $u,y$.

For two equations set
\[
 M_v(s=t,r=w):=\bigl(C(v,s,r)=C(v,t,w)\bigr).
\]
Extend by aligned fresh paired binders:
\[
\begin{aligned}
 M_v(D_{x,y}A,D_{x,y}B)&:=D_{x,y}M_v(A,B),\\
 M_v(E,D_{x,y}B)&:=D_{x,y}M_v(G_yE,B),\\
 M_v(D_{x,y}A,E)&:=D_{x,y}M_v(A,G_yE),
\end{aligned}
\]
where $E$ is an equation.  Each recursive call reduces maximum binder
depth.  The endpoint equivalences are
$M_0(A,B)\dashv\vdash A$ and $M_{Sv}(A,B)\dashv\vdash B$.

Put $\delta(v)=C(v,1,0)$ and $p(z)=I(z,0;u.\delta(u))$.  Only
$p(1)=1$ and $p(2)=0$ are needed.  The full inverse logical translation is
\[
\begin{aligned}
 A\Rightarrow_*B&:=D_{u,z}M_z(A,B),\\
 A\wedge_*B&:=D_{u,z}G_zM_{p(z)}(A,B),\\
 \exists_*x\,A(x)&:=D_{x,y}G_yA(x),\\
 A\vee_*B&:=\exists_*x\,M_x(A,B),\\
 \forall_*x\,A(x)&:=A(0)\wedge_*D_{u,z}G_zA(z),\\
 \bot_*&:=(0=1).
\end{aligned}
\]
All newly bound variables are fresh, and $u$ is unused.

Case introduction and the selector endpoints prove implication
introduction; elimination opens its clause and specializes at $1$.
The gate proves existential introduction and elimination.  Disjunction
uses witnesses $0$ and $1$, while its elimination uses the derived
numerical case rule.  Conjunction projections specialize at $1$ and
$2$; introduction uses numerical cases for the selector.  The first
universal conjunct supplies zero and the second supplies successors,
so coverage proves universal elimination.  Finally $0=1$ gives
$s=C(0,s,t)=C(1,s,t)=t$ for arbitrary terms; case introduction then
extends explosion to every formula.  These are the inherited logical
macro proofs, applicable because all case introductions and all
contextual equation inductions have been recovered in the new calculus.

The essential native formula composites are internally interderivable:
\[
 D_{x,y}B\dashv\vdash
 \exists_*x\bigl(B(x,0)\Rightarrow_*\forall_*y\,B(x,y)\bigr),
\]
\[
 Q_yB\dashv\vdash B(0)\Rightarrow_*\forall_*y\,B(y).
\]
For the first equivalence, paired elimination exposes the conditional
universal at a fresh witness; the converse uses paired case introduction
at the existential witness.  The second follows by closure elimination
and derived case-closure introduction.  External structural induction
on finite formulas proves both complete logical round trips.

\subsection{The inherited arithmetic bridge, with both composites}
The ConciseIPRA arithmetic development proves the following lemmas~\cite{bridge,readme}
\emph{internally in $T_I$}, in this order: a scalar predecessor
$P0=0$, $P(Sn)=n$; a term reflecting actual equality; a characteristic
term for every quantifier-free formula; pairing with two projections;
and a compiler for every fixed primitive-recursive description.  These
are inherited results, not new consequences asserted solely from
finite tests.  The source proofs are the reference files
\texttt{scalar\_predecessor.tex}, \texttt{equality\_reflection.tex}, and
\texttt{primitive\_recursion.tex}, with their corresponding Lean
modules named in the pinned repository README.

Here is the exact compiler step.  Suppose pairing $\pi$ and projections
$L,H$ have already been established, and a target PR description is
\[
 f(0,\vec p)=g(\vec p),\qquad
 f(Sn,\vec p)=h(n,f(n,\vec p),\vec p).
\]
For translations $g^*,h^*$ of its earlier descriptions put
\[
 V(n,\vec p)=I\bigl(n,\pi(0,g^*(\vec p));
 u.\pi(SL(u),h^*(L(u),H(u),\vec p))\bigr),
 \qquad f^*(n,\vec p)=H(V(n,\vec p)).
\]
Equation induction proves $L(V(n,\vec p))=n$.  The two projection
identities then establish both required defining equations for $f^*$.
The compiler recurses externally on a fixed finite PR description.

Conversely, iPRA interprets each iterator by primitive recursion whose
update ignores the old counter.  It proves the computation equations,
coverage, separation, and every literal-equation induction instance.
In the other direction, the predecessor gives successor injectivity,
the zero test gives successor separation, and the characteristic terms
reduce quantifier-free induction to literal-equation induction.

For the numerical round trips, two recursive operations with the same
seed and provably equal updates are internally equal: induction on the
literal equation between their outputs proves this at every counter.
External induction on finite descriptions therefore shows that every
operation is provably equal to its double translation.  Both domains
are whole and both interpreted equalities are actual equality, so
$x=y$ is an internally definable isomorphism from each structure to its
double interpretation.  Formula induction, using the logical composites
above, gives both formula round trips.  Effective proof translations
preserve contradiction in both directions and hence give
equiconsistency; they make no assertion that either theory proves its
own consistency.

Consequently the new calculus is a faithful presentation of an ordinary
theory bi-interpretable and equiconsistent with full first-order iPRA.
The proof of this conclusion uses the new exact sequent equivalence and
the inherited arithmetic bridge.  It does not rely on a standard-model
test, a quotient equality, or an interpretation restricted to standard
numerals.
